\documentclass{article}

% Language setting
% Replace `english' with e.g. `spanish' to change the document language
\usepackage[english]{babel}

% Set page size and margins
% Replace `letterpaper' with `a4paper' for UK/EU standard size
\usepackage[letterpaper,top=2cm,bottom=2cm,left=3cm,right=3cm,marginparwidth=1.75cm]{geometry}

% Useful packages
\usepackage{amsmath}
\usepackage{amssymb}
\usepackage{graphicx}
\usepackage[colorlinks=true, allcolors=blue]{hyperref}
\newcommand*\diff{\mathop{}\!\mathrm{d}}
\newcommand*\sndff{\mathbb{I}}

\title{Theorema Egregium}
\author{Michael Qian}

\begin{document}
\maketitle

\begin{abstract}
  Following Gauss's footsteps, we develop a method for quantifying the curvature of 2-dimensional surfaces, then show that this notion of curvature can be used as a ``compatibility condition" when determining whether or not one surface can serve as a faithful cartographic map of another--i.e., one which which preserves angles, distances (up to a constant scale factor $k$), and areas (up to a scale factor $k^2$). We then apply this compatibility condition to the analysis of two real-world problems: (1) mapping the surface of the Earth and (2) optimal pizza-wielding technique.
\end{abstract}

\section{Basic Definitions}

A \textbf{parametric surface} is defined by a smooth function $\vec{r} : D \to \mathbb{R}^3$ where $D$ is an open, connected subset of $\mathbb{R}^2$. $D$ acts the as the coordinate space used to parameterize the surface, while $\vec{r}$ maps each coordinate pair $(u_1, u_2) \in D$ to the point in $\mathbb{R}^3$ representing the physical location of the point on the surface described by those coordinates. The surface $S$ is then formed as the image of the entire coordinate space under the map $\vec{r}$:

\begin{equation*}
  S := \vec{r}(D) := \{\vec{r}(u_1, v_1) \mid (u_1, u_2) \in D\}
\end{equation*}

We will additionally assume that $\vec{r}$ is \textbf{regular}--the partial derivatives of $\vec{r}$ with respect to its first and second arguments (resp. $\partial_1 \vec{r}$ and $\partial_2 \vec{r}$) are linearly independent whenever they are evaluated at the same point $(u_1, u_2) \in D$, for all points in $D$. This guarantees that, at every point on the surface, varying one coordinate results in motion which cannot be achieved by varying the other coordinate. The smoothness and regularity of $\vec{r}$ together justify the intuition that $S$ is a 2-dimensional object embedded in 3-dimensional space.

Given a pair of smooth functions $u_1, u_2 : I \to D$ where $I$ is an open, non-empty interval of $\mathbb{R}$, we can define a curve on $S$ given by:

\begin{align*}
  \vec{\gamma} : I &\to \mathbb{R}^3 \\
  t &\mapsto \vec{r}(u_1(t), u_2(t))
\end{align*}

Each function $u_i(t)$ can be interpreted as indicating the value of the i\textsuperscript{th} coordinate at time $t$. The function $\vec{\gamma}$ then uses these coordinate functions to define a curve in $\mathbb{R}^3$ which is contained in $S$. The velocity of such a curve at time $t$ is given by:
\begin{equation*}
  \dot{\vec{\gamma}}(t) = \dot{u}_1(t) [\partial_1 \vec{r}](u_1(t), u_2(t)) + \dot{u}_2(t) [\partial_2 \vec{r}](u_1(t), u_2(t))
\end{equation*}

Where $\dot{u}_i$ is the $i$\textsuperscript{th} component of the \textbf{coordinate velocity}--the rate of change of the coordinate tuple $(u_1(t), u_2(t))$ with respect to $t$. Since the velocity $\dot{\vec{\gamma}}(t)$ is tangent to the curve at the point $\vec{\gamma}(t)$, and $\vec{\gamma}$ is fully contained in $S$ by construction, $\dot{\vec{\gamma}}(t)$ is therefore also tangent to the surface $S$ at the point $\vec{\gamma}(t)$. This motivates the definition of the the \textbf{tangent plane} $T_{\vec{r}(u_1, u_2)}S$ at the point $\vec{r}(u_1, u_2)$ for each $(u_1, u_2) \in D$ as the set of vectors given by:

\begin{equation*}
  T_{\vec{r}(u_1, u_2)}S := \{\alpha [\partial_1 \vec{r}](u_1, u_2) + \beta [\partial_2 \vec{r}](u_1, u_2) \mid (\alpha, \beta) \in \mathbb{R}^2\}
\end{equation*}

This plane can be understood as the set of all velocity vectors associated with all possible curves on $S$ passing through the point $\vec{r}(u_1, u_2)$. It is evident that this is a 2-dimensional subspace of $\mathbb{R}^3$ since it is defined as the linear span of the vectors $[\partial_1 \vec{r}](u_1, u_2)$ and $[\partial_2 \vec{r}](u_1, u_2)$.

Additionally, we define the \textbf{surface normal} $\vec{N}(u_1, u_2)$ and \textbf{unit surface normal} $\hat{n}(u_1, u_2)$ vectors at $\vec{r}(u_1, u_2)$ as:

\begin{align*}
  \vec{N}(u_1, u_2) &:= [\partial_1 \vec{r}](u_1, u_2) \times [\partial_2 \vec{r}](u_1, u_2) \\
  \hat{n}(u_1, u_2) &:= \frac{\vec{N}(u_1, u_2)}{|\vec{N}(u_1, u_2)|}
\end{align*}

These vectors stand orthogonal to $T_{\vec{r}(u_1, u_2)}S$ (and therefore also $S$ itself) at the point $\vec{r}(u_1, u_2)$ and are guaranteed to be non-zero everywhere on $S$ due to the regularity of $\vec{r}$. Note that the set $\{[\partial_1 \vec{r}](u_1, u_2), [\partial_2 \vec{r}](u_1, u_2), \hat{n}(u_1, u_2)\}$ is a linearly independent set of three vectors in $\mathbb{R}^3$, so any vector in $\mathbb{R}^3$ can be decomposed as a unique linear combination of them. The same is true if $\vec{N}(u_1, u_2)$ is substituted for $\hat{n}(u_1, u_2)$.

The unit surface normal $\hat{n}$ can additionally be interpreted as a map from $D$ to the unit sphere $\mathbb{S}^2$ (\textit{not} the Cartesian product of the surface $S$ with itself) due to its output always being unit-length. Under this interpretation, it is known as the \textbf{Gauss map}.

\section{Notation}

The functions ${\vec{r}}$, ${\vec{N}}$, ${\hat{n}}$, and all of their partial derivatives are implicitly evaluated at the same point $(u_1, u_2) \in D$ whenever they appear in an expression without being given explicit arguments, so that they are shorthands for the $\mathbb{R}^3$ vectors resulting from those evaluations. The same is true for all quantities constructed in terms of them. Any statement made about these quantities implicitly carries the qualifier ``for all $(u_1, u_2) \in D$."

\section{The First Fundamental Form}

Suppose we are given two tangent vectors $\vec{v}_1, \vec{v}_2 \in T_{\vec{r}}S$ in terms of their components along $\partial_1 \vec{r}$ and $\partial_2 \vec{r}$. What is the dot product of $\vec{v}_1$ and $\vec{v}_2$, considered as elements of $\mathbb{R}^3$?

Let:
\begin{align*}
  \vec{v}_1 = \alpha_1 \partial_1 \vec{r} + \beta_1 \partial_2 \vec{r} \\
  \vec{v}_2 = \alpha_2 \partial_1 \vec{r} + \beta_2 \partial_2 \vec{r}
\end{align*}

Then:
\begin{align*}
  \vec{v}_1 \cdot \vec{v}_2 &= (\alpha_1 \partial_1 \vec{r} + \beta_1 \partial_2 \vec{r}) \cdot (\alpha_2 \partial_1 \vec{r} + \beta_2 \partial_2 \vec{r}) \\
  &= \alpha_1 \alpha_2 (\partial_1 \vec{r} \cdot \partial_1 \vec{r}) + (\alpha_1 \beta_2 + \beta_1 \alpha_2) (\partial_1 \vec{r} \cdot \partial_2 \vec{r}) + \beta_1 \beta_2 (\partial_2 \vec{r} \cdot \partial_2 \vec{r})
\end{align*}

We now define the following quantities based on dot products formed using $\partial_1 {\vec{r}}$ and $\partial_2 {\vec{r}}$ and collect them into a matrix called the \textbf{first fundamental form}, denoted by the roman numeral $\mathrm{I}$:

\begin{align}
  E &= \partial_1 \vec{r} \cdot \partial_1 \vec{r} \\
  F &= \partial_1 \vec{r} \cdot \partial_2 \vec{r} \\
  G &= \partial_2 \vec{r} \cdot \partial_2 \vec{r}
\end{align}

\begin{equation}
  \mathrm{I} =
  \begin{bmatrix} E & F \\ F & G
  \end{bmatrix}
\end{equation}

This allows the preceding dot product to be written as:

\begin{align*}
  \vec{v}_1 \cdot \vec{v}_2 &= \alpha_1 \alpha_2 E + (\alpha_1 \beta_2 + \beta_1 \alpha_2) F + \beta_1 \beta_2 G \\
  &=
  \begin{bmatrix} \alpha_1 & \beta_1
  \end{bmatrix}
  \mathrm{I}
  \begin{bmatrix} \alpha_2 \\ \beta_2
  \end{bmatrix} \\
  &=
  \begin{bmatrix} \alpha_2 & \beta_2
  \end{bmatrix}
  \mathrm{I}
  \begin{bmatrix} \alpha_1 \\ \beta_1
  \end{bmatrix}
\end{align*}

The first fundamental form therefore allows us to compute the physical-space ($\mathbb{R}^3$) dot product of the tangent vectors $\vec{v}_1, \vec{v}_2 \in T_{\vec{r}}S$ in terms of their respective coordinate-space representations $(\alpha_1, \beta_1), (\alpha_2, \beta_2) \in \mathbb{R}^2$. It also aids in the computation of areas--given a patch of coordinate space $E \subseteq D$, the area of the corresponding patch of $S$ covered by $\vec{r}(E)$ is:

\begin{align*}
  A_{\vec{r}(E)} &= \int_E (\partial_1 \vec{r} \times \partial_2 \vec{r}) \cdot \hat{n} \diff{u_1} \diff{u_2} \\
  &= \int_E |\partial_1 \vec{r} \times \partial_2 \vec{r}| \diff{u_1} \diff{u_2} \\
  &= \int_E \sqrt{(\partial_1 \vec{r} \cdot \partial_1 \vec{r})(\partial_2 \vec{r} \cdot \partial_2 \vec{r}) - (\partial_1 \vec{r} \cdot \partial_2 \vec{r})^2} \diff{u_1} \diff{u_2} \\
  &= \int_E \sqrt{EG - F^2} \diff{u_1} \diff{u_2} \\
  &= \int_E \sqrt{\det(\mathrm{I})} \diff{u_1} \diff{u_2}
\end{align*}

Note that the assumed linear independence of $\partial_1 \vec{r}$ and $\partial_2 \vec{r}$ guarantees that $\det(\mathrm{I}) > 0$.

While our construction of the first fundamental form makes explicit reference to the way $S$ is embedded in $\mathbb{R}^3$ (since it is defined in terms of $\vec{r}$), one can alternatively consider defining an abstract surface by specifying the first fundamental form directly--i.e. specifying the coefficients $E$, $F$, and $G$ without reference to a higher-dimensional ambient space. By computing geometric quantities such as angles, lengths, and areas using only the first fundamental form, one can study the so-called \textit{intrinsic geometry} of the surface, which depends only on the internal properties of the surface itself. This is the approach taken by modern differential geometry and general relativity. Our own investigation of curvature in following sections will involve the construction and usage of quantiites which are ``intrinsic" in this sense.

\section{The Second Fundamental Form}

\section{The Gaussian Curvature}

\section{Isometries}

\section{Isometries Preserve Gaussian Curvature}

\section{Appendices}

\section{Shelf}

\begin{align*}
  \diff A_{\vec{n}(E)} &= (\partial_1 \hat{n} \times \partial_2 \hat{n}) \cdot \hat{n} \\
  &=  (\partial_1 (\vec{N} / |\vec{N}|) \times \partial_2 (\vec{N} / |\vec{N}|)) \cdot \hat{n} \\
  &= \frac{1}{|\vec{N}|^2} [(|\vec{N}| \partial_1 \vec{N} - (\partial_1 |\vec{N}|)\vec{N}) \times (|\vec{N}| \partial_2 \vec{N} - (\partial_2|\vec{N}|)\vec{N})] \cdot \hat{n} \\
  &= \frac{1}{|\vec{N}|^2} [|\vec{N}|^2(\partial_1 \vec{N} \times \partial_2 \vec{N}) - |\vec{N}|(\partial_1 |\vec{N}|)(\vec{N} \times \partial_2 \vec{N}) - |\vec{N}| (\partial_2 |\vec{N}|)(\partial_1 \vec{N} \times \vec{N})] \cdot \hat{n} \\
  &= (\partial_1 \vec{N} \times \partial_2 \vec{N}) \cdot \hat{n}
\end{align*}

\section{What exactly do we mean by ``curvature"?}

We begin our exploration of surface curvature by investigating, for a given patch of coordinate space $E \subseteq D$, the relationship between the signed area $A_{\vec{r}(E)}$ on $S$ covered by $\vec{r}(E)$ and the signed area $A_{\hat{n}(E)}$ on $\mathbb{S}^2$ covered by $\hat{n}(E)$:

\begin{align*}
  A_{\vec{r}(E)} &= \int_E (\partial_1 \vec{r} \times \partial_2 \vec{r}) \cdot \hat{n} \diff{u_1} \diff{u_2} \\
  A_{\hat{n}(E)} &= \int_E (\partial_1 \hat{n} \times \partial_2 \hat{n}) \cdot \hat{n} \diff{u_1} \diff{u_2}
\end{align*}

The motivation for doing so is this: $\mathbb{S}^2$ effectively represents the ``space of directions" one can face in $\mathbb{R}^3$. The area $A_{\hat{n}(E)}$ then represents the portion of this ``space of directions" covered by the unit surface normals of $\vec{r}(E)$. This suggests a definition of curvature based on ``amount of direction space covered per unit surface area," which directs our attention to the relationship between the integrands.

First, we analyze the partial derivatives $\partial_i\hat{n}$:

\begin{align*}
  \partial_i \hat{n} &= \partial_i \left( \frac{\vec{N}}{|\vec{N}|} \right) \\
  &= \frac{1}{|\vec{N}|^2}(|\vec{N}| \partial_i \vec{N} - (\partial_i |\vec{N}|) \vec{N}) \\
&= \frac{1}{|\vec{N}|^2} (|\vec{N}| \partial_i \vec{N} - (\partial_i ((\vec{N} \cdot \vec{N})^{\frac{1}{2}}))\vec{N})) \\
&= \frac{1}{|\vec{N}|^2} \left(|\vec{N}| \partial_i \vec{N} - \left( \frac{1}{2} (\vec{N} \cdot \vec{N})^{-\frac{1}{2}} (2 (\partial_i \vec{N}) \cdot \vec{N}) \right) \vec{N} \right) \\
&= \frac{1}{|\vec{N}|^2} \left(|\vec{N}| \partial_i \vec{N} - \left( \frac{1}{|\vec{N}|} (\partial_i \vec{N}) \cdot \vec{N} \right) \vec{N} \right) \\
&= \frac{1}{|\vec{N}|} (\partial_i \vec{N} - ((\partial_i \vec{N}) \cdot \hat{n}) \hat{n})
\end{align*}

The expression $(\partial_i \vec{N} - ((\partial_i \vec{N}) \cdot \hat{n}) \hat{n})$ subtracts from $\partial_i \vec{N}$ the component of itself which is parallel to $\hat{n}$, leaving only the orthogonal component which we denote by $(\partial_i \vec{N})_{\perp \hat{n}}$. $(\partial_i \vec{N})_{\perp \hat{n}}$ is therefore an element of the tangent plane $T_{\vec{r}}S$, so there exist coefficients $\alpha$ and $\beta$ such that:

\begin{align*}
\partial_i \hat{n} = \frac{(\partial_i \vec{N})_{\perp \hat{n}}}{|\vec{N}|} = \alpha \partial_1 \vec{r} + \beta \partial_2 \vec{r}
\end{align*}

To find these coefficients, we first define the following values based on dot products formed using $\partial_1 {\vec{r}}$ and $\partial_2 {\vec{r}}$:

\begin{align}
E &= \partial_1{\vec{r}} \cdot \partial_1{\vec{r}} \\
F &= \partial_1{\vec{r}} \cdot \partial_2{\vec{r}} \\
G &= \partial_2{\vec{r}} \cdot \partial_2{\vec{r}}
\end{align}

We then have:

{
\newcommand{\ndrone}{\partial_i \hat{n} \cdot \partial_1{\vec{r}}}
\newcommand{\ndrtwo}{\partial_i \hat{n} \cdot \partial_2{\vec{r}}}

\begin{align}
  \label{eq4}
  \ndrone &= \alpha E + \beta F \\
  \label{eq5}
  \ndrtwo &= \alpha F + \beta G
\end{align}

We can cancel all terms involving $\beta$ by multiplying Eq. \eqref{eq4} by $G$, multiplying Eq. \eqref{eq5} by $F$, then subtracting the two. Rearranging to isolate $\alpha$ gives:

\begin{equation}
  \label{eq6}
  \alpha = \frac{(\ndrone)G - (\ndrtwo)F}{EG - F^2}
\end{equation}

A similar manipulation can be used to solve for $\beta$, giving:

\begin{equation}
  \label{eq7}
  \beta = \frac{(\ndrtwo)E - (\ndrone)F}{EG - F^2}
\end{equation}

}

Note that the assumed linear independence of $\partial_1 \vec{r}$ and $\partial_2 \vec{r}$ guarantees that:

\begin{equation*}
EG - F^2 = (\partial_1 \vec{r} \cdot \partial_1 \vec{r})(\partial_2 \vec{r} \cdot \partial_2 \vec{r}) - (\partial_1 \vec{r} \cdot \partial_2 \vec{r})^2 > 0
\end{equation*}

\end{document}
